%----------------------------------------------------------------------
\section{Introduction}\label{sec:intro}
%----------------------------------------------------------------------

The companion manuscript \emph{Discrimination as Foundation}
(henceforth \textbf{[DAF]}) establishes a constructive foundation
for set theory in which the primitive notion is discrimination, valued
in the exterior algebra $\bigwedge(\GF(2)^n)$, with a four-cell
membership profile $(\tau, \sigma) \in \GF(2)^2$ and a monotonically
growing regime $\kappa$. The framework reaches a definite resting
point: set-theoretic consequences (extensionality, Russell exclusion,
foundation), categorical consequences (adjunctions, Yoneda, topos
structure), and homotopical consequences (chain complexes, univalence)
are all derivable from the algebra plus three axioms (profile
linearity \textsf{P}, evaluation linearity \textsf{E}, operationalist
principle \textsf{O}).

That resting point is mathematical, not operational. [DAF]
establishes \emph{what is computed}; it is silent about
\emph{how to compute it efficiently on modern hardware}. The present
paper fills that gap. The concrete deliverable is a compiler
(``planner'') that takes a DAG of CSTZ atomic operations and
produces one or more executable plans, each a concrete schedule of
dispatches to a heterogeneous computing substrate (multi-tier memory,
SIMT/MIMD warps, register files of finite depth). The compiler is
architecturally small but mathematically interesting: we argue that
it is most naturally expressed as a weighted attribute grammar whose
terminals are the CSTZ atoms, whose productions are the admissible
scheduling transformations, and whose attributes are cost vectors
over four resources. The compiler's inner loop is then not an
ad hoc dispatch table but a chart parser.

\subsection{The execution problem}

Let $W$ be a workload: a finite DAG of atomic $\GF(2)$ operations
over exterior elements. In [DAF], $W$ arises as the sequence of
discriminator applications, residue computations, $\kappa$-evolutions,
and sheafification steps that the inference engine executes on a
bitstream. Each atom is mathematically trivial (a handful of XORs
and ANDs); there are many of them, and they must share a small
hardware substrate.

Three structural observations drive what follows.

\begin{enumerate}[label=(\roman*)]
  \item \textbf{Encoding is a phase.} Every exterior element has at
    least three equivalent representations: the coefficient basis
    $\Lambda$, the Reed--Muller syndrome basis RM (related to
    $\Lambda$ by the zeta transform $\zeta$, which is
    self-inverse over $\GF(2)$), and a grade-truncated or
    Hamming-compressed form for objects lying in a proper subcode.
    Operators have different cost in different bases: $\partial$
    collapses the RM grade budget monotonically, $\wedge$ is
    pointwise multiply after ranked zeta, $\mathrm{kappa\_equiv}$
    is a single syndrome check in the RM basis. The planner's
    primary job is to \emph{assign encodings to subplans}.

  \item \textbf{Waste is forced batch coarseness.} Hardware tile
    sizes (e.g.\ 4, 8, 16 per axis on tensor-core pipelines) do not
    in general divide the natural CSTZ tile sizes (2 for perspective
    rotations, 3 for Fano syndromes, $n$ for general $\GF(2)^n$
    operators). Under LCM packing, the utilization of a
    single hardware block of size $M \times M$ by CSTZ tiles of
    size $d \times d$ is 100\% regardless of $d$; what differs
    between tile sizes is the minimum batch size $(L/d)^2$ (where
    $L = \lcm(d,M)$) that saturates the block. The ``wasted'' lanes
    under naive padding are in fact a budget for additional work:
    either speculative execution of alternate plans, or
    triple-modular-redundant execution of the same plan for fault
    tolerance.

  \item \textbf{Multi-resource optimization is multi-objective.}
    Four budgets compete: VRAM (device-global memory), shared
    memory (per-SM staging), register file (per-thread live values),
    and wall-clock latency. Fusion trades VRAM bandwidth against
    SMEM and register pressure; time-sharing trades latency against
    register pressure. The planner cannot collapse these four into
    a single scalar cost without loss of information, but it can
    enforce them jointly via Lagrangian duality over the resource
    vector. The dual prices give a direct diagnostic of which
    resource is currently binding, architecture-by-architecture.
\end{enumerate}

These observations converge on a design: a weighted attribute
grammar whose chart parser is an Earley variant with cost-bounded
completion. The grammar's productions are scheduling transformations
(fuse linear operators, time-share a register, re-encode a subplan);
the attributes are cost vectors; the parser's completer rejects
chart items whose aggregate cost would exceed the budget. Optimality
under hard budgets is the Lagrangian dual of chart parsing; the
output portfolio is the Pareto-nondominated frontier of that parse.

\subsection{The successor axiom}

[DAF]'s operationalist axiom states that a distinction that cannot
be witnessed by any discriminator under any $\kappa$-evolution is
not a valid distinction. We add an execution-level analogue, labeled
\textsf{O}$_{\mathrm{exec}}$ to distinguish it from the foundational
axiom:

\begin{definition}[Execution-operationalist axiom]
  \label{def:exec-operationalist}
  Let $P_1$ and $P_2$ be two plans for the same workload $W$, both
  valid under the planner grammar, and let $c(P_i) \in
  \mathbb{N}^4$ denote the cost vector of $P_i$
  under a stated hardware profile. If for every resource $r$ in the
  budget vector $B$ we have $c(P_1)_r \leq B_r$ and $c(P_2)_r \leq
  B_r$, and if no subsequent measurement (latency, bytes, flops,
  register spills) can distinguish $P_1$ from $P_2$ on the target
  hardware, then $P_1$ and $P_2$ are \emph{execution-equivalent} on
  that hardware. The planner is free to emit either, or a
  lane-packed mixture of both.
\end{definition}

The execution-operationalist axiom is to plans what [DAF]'s
operationalist axiom is to discriminations: it declares the
execution-equivalence relation to be the coarsest relation
consistent with observable cost behavior on a stated substrate.
It licenses the portfolio output---if two plans are cost-indistinguishable,
the planner need not choose---and it is the formal basis for the
ambiguous-parse packing of Section \ref{sec:portfolios}.

\subsection{Reading guide: how to verify this paper's claims}
\label{sec:reading-guide}

The paper contains 42 formal objects across nine sections.
Before encountering them, the reader should know three
things: where each claim is verified, what is not claimed,
and how much of the paper depends on which axioms.

\medskip
\noindent\textbf{Axiom-class distribution.}
The paper rests on commitments of three kinds. The first
two are inherited from [DAF] and appear with the same
labels: \textsf{A} (algebra of $\GF(2)$ and $\bigwedge V$),
\textsf{P} (profile linearity), \textsf{E} (evaluation
linearity), \textsf{O} (operationalist). The third is
new to this paper, reflecting its hardware-grounding:

\begin{description}[nosep,leftmargin=1.5cm]
  \item[\textsf{H}] \emph{Hardware profile}: the claim
    depends on a stated cost model assigning $(\VRAM,
    \SMEM, \REG, \LAT)$ values to each production. The
    profile is a per-architecture calibration; the claim
    is valid relative to any profile that satisfies the
    model's monotonicity axioms (Section \ref{sec:resource-vector}).
  \item[\textsf{R}] \emph{Reversibility}: the claim
    depends on GF(2)-linearity of specific atoms for the
    XOR-bridge mechanic to apply. Non-linear atoms
    (popcount-reduce, comparison) break \textsf{R} and
    require an explicit register materialization.
  \item[\textsf{F}] \emph{Fusion safety}: the claim
    depends on lane-disjointness of the masks assigned
    to fused atoms. Non-disjoint masks invalidate the
    block-diagonal decomposition.
  \item[\textsf{O}$_{\mathrm{exec}}$] The
    execution-operationalist axiom
    (Definition \ref{def:exec-operationalist}).
    A plan-equivalence relation rather than a
    set-equivalence relation.
\end{description}

Of the 42 formal objects: 17 (40\%) depend only on
\textsf{A}. 14 additionally require \textsf{H}. 6 require
\textsf{AHR} (reversibility). 3 require \textsf{AHF} (fusion
safety). 2 require \textsf{AHO}$_{\mathrm{exec}}$ (the new
axiom, touching only portfolio extraction and Pareto
packing).

\medskip
\noindent\textbf{Verification map.}
Each claim's proof lives in the spine (\S\S3--8); concrete
hardware-grounded witnesses live in Appendix \ref{app:kernels}
(Pallas kernel sketches) and Appendix \ref{app:examples}
(worked examples). Each theorem carries a footnote stating:
axiom class, dependency depth, verification pointer, warrant,
qualifier, finite-resource validity, and novelty status.

\begin{center}
\scriptsize
\begin{tabular}{l|c|c|l}
\textbf{Claim} & \textbf{Class} & \textbf{Status} &
  \textbf{Verification} \\
\hline
Zeta involutivity (\ref{prop:zeta-involution}) & \textsf{A} & \cmark &
  Double-sum, $1+1=0$ \\
Encoding commutation (\ref{prop:encoding-commute}) & \textsf{A} & \cmark &
  $\zeta$ $\GF(2)$-linear; commutes with $\oplus$ \\
LCM density (\ref{prop:lcm-density}) & \textsf{AH} & \cmark &
  Counting; $(M/d)^2 d^2 = M^2$ \\
XOR-bridge correctness (\ref{prop:xor-bridge}) & \textsf{AHR} & \cmark &
  App.\ \ref{app:b}; linearity of $L$ \\
Fused-apply correctness (\ref{prop:fused-correctness}) & \textsf{AHF} & \cmark &
  App.\ \ref{app:b}; block-diagonal decomposition \\
Time-share commutativity (\ref{prop:ts-commute}) & \textsf{AHR} & \cmark &
  Lane-disjoint masks + linearity \\
Grammar soundness (\ref{thm:soundness}) & \textsf{A} & \cmark &
  Induction on derivations \\
Cost-bounded Earley completeness (\ref{thm:earley-complete}) & \textsf{AH} & \cmark &
  Admissible heuristic; App.\ \ref{app:b} \\
Planner complexity (\ref{thm:complexity}) & \textsf{AH} & \cmark &
  $O(N^3 \cdot |\mathrm{Pareto}|)$; std. Earley \\
Lagrangian achievability (\ref{thm:lagrangian}) & \textsf{AH} & \omark &
  LP relaxation; integrality gap bounded \\
Pareto portfolio extraction (\ref{prop:pareto}) & \textsf{AH} & \cmark &
  Dominance check on chart \\
Parse-packing soundness (\ref{thm:parse-packing}) & \textsf{AHO}$_{\mathrm{exec}}$ & \cmark &
  Mask disjointness + \textsf{O}$_{\mathrm{exec}}$ \\
Plan-equivalence transitivity (\ref{prop:plan-trans}) & \textsf{AO}$_{\mathrm{exec}}$ & \cmark &
  Direct from Definition \ref{def:exec-operationalist} \\
Register-time-space bound (\ref{thm:reg-bound}) & \textsf{AH} & \cmark &
  Interval graph; App.\ \ref{app:b} \\
Incremental replanning (\ref{prop:incremental}) & \textsf{AH} & \cmark &
  Earley chart locality \\
\end{tabular}
\end{center}

\noindent Legend: \cmark\ = verified,
\gmark\ = gap (specific issue identified),
\omark\ = open-acknowledged.

\medskip
\noindent\textbf{What is not claimed (open questions).}

\begin{description}[nosep,leftmargin=0.8cm]
  \item[X1] \emph{Optimal calibration from first principles.}
    The cost model presupposes a hardware profile; we do
    not derive those numbers analytically. A principled
    calibration from published SM specifications is
    plausible but not attempted. \textbf{Open.}
  \item[X2] \emph{Non-linear atom fusion.}
    The XOR-bridge mechanic requires GF(2)-linearity.
    We conjecture (Conjecture \ref{conj:nonlinear}) that
    certain non-linear atoms (popcount, comparison)
    admit a weaker ``semi-reversibility'' via a paired
    inverse atom, but we do not prove it.
  \item[X3] \emph{Integrality gap for Lagrangian LP.}
    We show an LP relaxation is tractable, but the
    gap between LP-optimal and IP-optimal plan cost
    is bounded only empirically. A tight worst-case
    bound is \textbf{open-acknowledged}.
  \item[X4] \emph{Dynamic shapes under JIT.}
    The cost model presupposes static $n$ (exterior
    dimension). Dynamic $n$ is expressible as a
    family of plans indexed by $n$, but the family's
    joint optimization is outside this paper's scope.
\end{description}

\medskip
\noindent\textbf{Gap priorities} (from downstream
impact on the DAF→planner→kernel pipeline).

\emph{Tier 1 --- load-bearing}
(closing these unblocks the most downstream results):

\begin{description}[nosep,leftmargin=0.8cm]
  \item[S1] Calibration profile for a reference GPU
    (A100 or H100). \textbf{Open.} Unblocks all
    \textsf{H}-class claims with concrete numerics.
  \item[S2] Implementation of the Earley chart parser
    against the atom inventory of Section \ref{sec:grammar}.
    \textbf{Open.} Unblocks integration with the
    existing Pallas kernel layer.
\end{description}

\emph{Tier 2 --- structural}:

\begin{description}[nosep,leftmargin=0.8cm]
  \item[S3] Grammar well-formedness check (no
    left-recursion on non-linear atoms; admissible
    heuristic is monotone). \textbf{Partial.}
    Sufficient conditions stated; grammar-level
    decidability deferred.
  \item[S4] Property-based equivalence test against
    a naive reference implementation.
    \textbf{Outlined; not executed.}
\end{description}

\subsection{What changes from [DAF]}

The successor axiom \textsf{O}$_{\mathrm{exec}}$ is the only
genuinely new foundational commitment. Everything else is a
lifting of [DAF]'s algebraic machinery into an execution model:

\begin{itemize}[nosep]
  \item [DAF]'s four-cell structure becomes the atom encoding
    tag (coefficient/syndrome/Hamming/truncated).
  \item [DAF]'s $\kappa$-evolution becomes the chart parser's
    forward pass: each atom in the workload DAG extends the
    chart monotonically, in strict analogy to how each
    $\kappa$-evolution extends the regime.
  \item [DAF]'s Yoneda-style identification of objects by
    morphism profiles becomes the equivalence relation on
    plans: plans are identified by their observable cost
    signature (\textsf{O}$_{\mathrm{exec}}$).
  \item [DAF]'s Fano-plane and Reed--Muller structure at
    Section 9 appears here as the native tile size for
    Hamming-compressed syndrome operations
    (Proposition \ref{prop:lcm-density}).
\end{itemize}

The paper's novelty is not in the mathematical content of any
individual lifting but in the synthesis: treating the
compiler's internal data structure as a chart in the same
sense that [DAF] treats the discrimination kernel's internal
state as a sheaf, with cost vectors playing the role of
morphism profiles.

\subsection{Roadmap}

Section \ref{sec:prelim} fixes notation, the encoding lattice,
the resource vector, and the terminology concordance with
standard compiler and parser theory. Section \ref{sec:grammar}
presents the atom inventory and base productions. Sections
\ref{sec:fusion} and \ref{sec:time-sharing} develop the two
compound productions---spatial (MIMD lane fusion) and
temporal (XOR-bridge register time-sharing)---and prove
their correctness. Section \ref{sec:parsing} specifies the
Earley chart parser with cost-bounded completion and gives
the complexity bound. Section \ref{sec:lagrangian} derives
the dual formulation and the dual-price diagnostic.
Section \ref{sec:portfolios} extracts the Pareto frontier
and proves soundness of ambiguous-parse packing. Section
\ref{sec:discussion} summarizes, relates to [DAF], and
enumerates limitations. Appendix \ref{app:analogies}
collects structural analogies to classical compiler and
parser theory; Appendix \ref{app:b} gives detailed proofs;
Appendix \ref{app:kernels} sketches Pallas kernels;
Appendix \ref{app:meta} presents the closure-technique
matrix for this paper's proofs; Appendix \ref{app:examples}
works through a complete example end-to-end.
